\section{Énergie d’horizon et cellule de demi-action}
\label{rm:chap:horizonte-media-accion}

\subsection{Conventions géométriques de la région de Sitter}

Considérons un espace de Sitter quadridimensionnel de rayon aréal \(R>0\), de vitesse causale \(c\) et de constante gravitationnelle \(G\). Adoptons la convention
\begin{equation}
 H_R=\frac cR,
 \qquad
 \Lambda_R=\frac3{R^2},
 \qquad
 \ell_P^2=\frac{G\hbar}{c^3}.
 \label{rm:eq:horizon-ds-conventions}
\end{equation}
La densité d’énergie associée au terme cosmologique est
\begin{equation}
 \varepsilon_\Lambda
 :=\frac{\Lambda_Rc^4}{8\pi G}.
 \label{rm:eq:horizon-vacuum-density}
\end{equation}
La section spatiale plate de la région contient, jusqu’au rayon \(R\), le volume
\begin{equation}
 V_R=\frac{4\pi}{3}R^3.
 \label{rm:eq:horizon-flat-volume}
\end{equation}
Ces conventions fixent positivement l’énergie quasi locale utilisée dans ce chapitre. Le choix d’un autre signe pour la charge de Killing de Sitter relève d’une construction différente et doit être traduit avant toute comparaison des formules.

\subsection{Énergie du vide et énergie de Misner--Sharp}

Dans une géométrie à symétrie sphérique, soit
\begin{equation}
 \dd s^2=g_{ab}(x)\dd x^a\dd x^b
 +\mathcal R(x)^2\dd\Omega_2^2.
 \label{rm:eq:horizon-spherical-metric}
\end{equation}
L’énergie de Misner--Sharp \cite{rm:MisnerSharp1964}, dans la convention positive adoptée, est définie par
\begin{equation}
 E_{\rm MS}(\mathcal R)
 :=\frac{c^4\mathcal R}{2G}
 \left(1-g^{ab}\nabla_a\mathcal R\nabla_b\mathcal R\right).
 \label{rm:eq:horizon-Misner-Sharp}
\end{equation}
Une sphère marginale satisfait
\begin{equation}
 g^{ab}\nabla_a\mathcal R\nabla_b\mathcal R=0.
 \label{rm:eq:horizon-marginal-sphere}
\end{equation}

\begin{theorem}[Confluence volumétrique et quasi locale]
\label{rm:thm:horizon-vacuum-MS}
À l’horizon de Sitter \(\mathcal R=R\), l’énergie intégrée du vide et l’énergie de Misner--Sharp coïncident exactement :
\begin{equation}
 \boxed{
 E_\Lambda
 :=\varepsilon_\Lambda V_R
 =E_{\rm MS}(R)
 =\frac{c^4R}{2G}.}
 \label{rm:eq:horizon-E-Lambda}
\end{equation}
\end{theorem}

\begin{proof}
Les équations \cref{rm:eq:horizon-ds-conventions,rm:eq:horizon-vacuum-density,rm:eq:horizon-flat-volume} donnent
\[
 \varepsilon_\Lambda V_R
 =\frac{3c^4}{8\pi GR^2}\frac{4\pi R^3}{3}
 =\frac{c^4R}{2G}.
\]
La condition marginale \cref{rm:eq:horizon-marginal-sphere} appliquée à \cref{rm:eq:horizon-Misner-Sharp} produit le même résultat.
\end{proof}

L’égalité réunit deux expressions géométriques. La première intègre une densité volumétrique ; la seconde évalue le défaut de piégeage de la sphère aréale. Toutes deux conservent le rayon \(R\) et la compliance gravitationnelle \(G/c^4\).

\subsection{Température, entropie et produit d’horizon}

La gravité de surface de Sitter a pour module \(c^2/R\). La température de Gibbons--Hawking \cite{rm:GibbonsHawking1977} et l’entropie de Bekenstein--Hawking \cite{rm:Bekenstein1973,rm:Hawking1975} sont
\begin{align}
 k_BT_H&=\frac{\hbar c}{2\pi R},
 \label{rm:eq:horizon-temperature}\\
 \frac{S_H}{k_B}
 &=\frac{A_H}{4\ell_P^2}
 =\frac{4\pi R^2}{4\ell_P^2}
 =\frac{\pi R^2}{\ell_P^2}.
 \label{rm:eq:horizon-entropy}
\end{align}

\begin{theorem}[Identité thermodynamique de Sitter]
\label{rm:thm:horizon-TS}
Pour tout rayon \(R\) de la famille déclarée,
\begin{equation}
 \boxed{
 T_HS_H=\frac{c^4R}{2G}=E_\Lambda.}
 \label{rm:eq:horizon-TS-E}
\end{equation}
\end{theorem}

\begin{proof}
La multiplication de \cref{rm:eq:horizon-temperature,rm:eq:horizon-entropy} donne
\[
 T_HS_H
 =\frac{\hbar c}{2\pi Rk_B}
  \frac{k_B\pi R^2}{\ell_P^2}
 =\frac{\hbar cR}{2\ell_P^2}.
\]
L’identité \(\ell_P^2=G\hbar/c^3\) transforme le dernier membre en \(c^4R/(2G)\), égal à \cref{rm:eq:horizon-E-Lambda}.
\end{proof}

Le facteur \(1/2\) possède une décomposition géométrique transparente. Le facteur aréal \(4\pi\), la normalisation holographique par quatre et la périodicité thermique \((2\pi)^{-1}\) satisfont
\begin{equation}
 \frac{4\pi}{4}\frac1{2\pi}=\frac12.
 \label{rm:eq:horizon-half-factor}
\end{equation}
La moitié est donc la contraction conjointe entre quantification aréale et périodicité euclidienne.

\subsection{Quantum complet de l’horizon et force de Planck}

Définissons la tension ou force de Planck et le quantum complet associé au rayon par
\begin{equation}
 \mathscr F_P:=\frac{c^4}{G},
 \qquad
 \boxed{
 \mathcal Q_H(R):=2T_HS_H
 =\frac{c^4R}{G}=\mathscr F_P\,R.}
 \label{rm:eq:horizon-full-quantum}
\end{equation}
L’énergie quasi locale occupe une demi-cellule du quantum complet :
\begin{equation}
 \boxed{E_\Lambda=T_HS_H=\frac{\mathcal Q_H(R)}2.}
 \label{rm:eq:horizon-half-quantum-general}
\end{equation}

\begin{proposition}[Compliance radiale]
\label{rm:prop:horizon-compliance}
La dérivée du quantum complet par rapport au rayon est la force de Planck, et la dérivée de l’énergie quasi locale en est la moitié :
\begin{equation}
 \frac{\dd\mathcal Q_H}{\dd R}=\mathscr F_P,
 \qquad
 \frac{\dd E_\Lambda}{\dd R}=\frac{\mathscr F_P}{2}.
 \label{rm:eq:horizon-radial-derivatives}
\end{equation}
La compliance gravitationnelle \(\mathscr C_P=G/c^4\) satisfait
\begin{equation}
 R=\mathscr C_P\mathcal Q_H.
 \label{rm:eq:horizon-compliance-law}
\end{equation}
\end{proposition}

\begin{proof}
Les trois identités sont des dérivations ou des inversions directes de \cref{rm:eq:horizon-full-quantum}.
\end{proof}

L’équation \cref{rm:eq:horizon-compliance-law} exprime la réponse radiale : le quantum complet applique une tension et la géométrie exprime un rayon au moyen de la compliance universelle. Le produit \(T_HS_H\) occupe la moitié quasi locale de cette réponse.

\subsection{Spécialisation à la fermeture chronologique de 108 états}

La fermeture chronologique de \cref{rm:ch:cosmo} détermine
\begin{equation}
 R_{108}=\ell_P,
 \qquad
 t_P=\frac{\ell_P}{c},
 \qquad
 T_{108}=2\pi t_P,
 \label{rm:eq:horizon-CT108-scales}
\end{equation}
et le quantum de Planck
\begin{equation}
 E_P=\frac{\hbar}{t_P}
 =k_B\Theta_{\rm clk}\log3.
 \label{rm:eq:horizon-Planck-trit}
\end{equation}

\begin{theorem}[Cellule CT108 de demi-action]
\label{rm:thm:horizon-half-action}
Dans la spécialisation \(R=R_{108}=\ell_P\),
\begin{align}
 \mathcal Q_H(\ell_P)
 &=\frac{c^4\ell_P}{G}=E_P,
 \label{rm:eq:horizon-QH-EP}\\
 \boxed{
 E_\Lambda=T_HS_H
 =\frac{E_P}{2}
 =\frac{k_B\Theta_{\rm clk}\log3}{2},}
 \label{rm:eq:horizon-CT108-half-energy}\\
 \boxed{E_\Lambda t_P=\frac{\hbar}{2},}
 \qquad
 \boxed{E_\Lambda T_{108}=\frac h2.}
 \label{rm:eq:horizon-half-actions}
\end{align}
\end{theorem}

\begin{proof}
La première égalité s’obtient en substituant \(\ell_P=\sqrt{G\hbar/c^3}\) et \(E_P=\sqrt{\hbar c^5/G}\). \cref{rm:eq:horizon-half-quantum-general} produit \(E_\Lambda=E_P/2\), et \cref{rm:eq:horizon-Planck-trit} fournit la carte informationnelle. Enfin,
\[
 E_\Lambda t_P=\frac{E_Pt_P}{2}=\frac\hbar2
\]
et, puisque \(T_{108}=2\pi t_P\) et \(h=2\pi\hbar\),
\[
 E_\Lambda T_{108}
 =\frac{E_P}{2}(2\pi t_P)=\pi\hbar=\frac h2.
\]
\end{proof}

La première égalité de \cref{rm:eq:horizon-half-actions} exprime une demi-action réduite sur le temps de Planck ; la seconde exprime une demi-action complète sur la période CT108. Toutes deux représentent une même cellule dans les cartes angulaire et circulaire.

La température et l’entropie sont également déterminées :
\begin{equation}
 k_BT_H=\frac{E_P}{2\pi},
 \qquad
 \boxed{\frac{S_H}{k_B}=\pi.}
 \label{rm:eq:horizon-CT108-TS}
\end{equation}
La température convertit la fréquence d’horizon en énergie ; l’entropie convertit l’aire de la sphère de Planck en information. Leur produit retrouve le même demi-quantum.

\subsection{Poids effectif de l’horizon et exponentielle hyperbolique}

Définissons le poids thermodynamique effectif
\begin{equation}
 \Omega_{\rm eff}(R):=\exp\!\left(\frac{S_H(R)}{k_B}\right).
 \label{rm:eq:horizon-effective-weight}
\end{equation}
Au rayon CT108,
\begin{equation}
 \boxed{\Omega_{\rm eff}(\ell_P)=\ee^\pi.}
 \label{rm:eq:horizon-effective-e-pi}
\end{equation}
Le poids \(\Omega_{\rm eff}\) est une grandeur positive définie par l’entropie. Sa valeur réelle exprime une multiplicité effective ; l’interpréter comme cardinal entier de microétats exige une discrétisation supplémentaire.

L’opérateur d’incidence aréale--volumétrique est
\begin{equation}
 F_{\rm av}=\begin{pmatrix}0&1\\1&1\end{pmatrix},
 \qquad
 J_\varphi:=\frac{2F_{\rm av}^2-3I}{\sqrt5}.
 \label{rm:eq:horizon-J-phi}
\end{equation}
Un calcul direct donne
\begin{equation}
 J_\varphi^2=I,
 \qquad
 \Spec(J_\varphi)=\{+1,-1\}.
 \label{rm:eq:horizon-J-involution}
\end{equation}
Le flot hyperbolique au temps de fermeture \(\pi\) est
\begin{equation}
 \mathcal E_\pi
 :=\ee^{\pi J_\varphi}
 =\cosh\pi\,I+\sinh\pi\,J_\varphi,
 \label{rm:eq:horizon-hyperbolic-flow}
\end{equation}
et possède le spectre
\begin{equation}
 \boxed{
 \Spec(\mathcal E_\pi)=\{\ee^\pi,\ee^{-\pi}\},
 \qquad
 \rho(\mathcal E_\pi)=\ee^\pi.}
 \label{rm:eq:horizon-hyperbolic-spectrum}
\end{equation}

\begin{corollary}[Confluence scalaire horizon--Perron]
\label{rm:cor:horizon-Perron-scalar}
À la fermeture CT108,
\begin{equation}
 \boxed{
 \frac{S_H}{k_B}
 =\log\Omega_{\rm eff}
 =\log\rho(\mathcal E_\pi)
 =\pi.}
 \label{rm:eq:horizon-Perron-log}
\end{equation}
\end{corollary}

\begin{proof}
Les équations \cref{rm:eq:horizon-CT108-TS,rm:eq:horizon-effective-e-pi,rm:eq:horizon-hyperbolic-spectrum} donnent les trois membres.
\end{proof}

L’identité scalaire est exacte et conserve deux généalogies distinctes. L’horizon produit \(\ee^\pi\) par exponentiation de l’entropie aréale ; la dynamique de Perron le produit comme valeur propre expansive d’une involution hyperbolique. L’égalité terminale fixe une condition aux limites pour l’entrelaceur et préserve la distinction des domaines.

\subsection{Cohérence du transducteur gravitationnel avec l’horizon de Schwarzschild}

Un horizon de Schwarzschild de même rayon aréal satisfait
\begin{equation}
 R_H=\frac{2GE_H}{c^4},
 \qquad
 E_H=\frac{c^4R_H}{2G}.
 \label{rm:eq:horizon-Schwarzschild-energy}
\end{equation}
L’entropie aréale coïncide avec \cref{rm:eq:horizon-entropy}, tandis que la température est
\begin{equation}
 k_BT_{\rm Sch}=\frac{\hbar c}{4\pi R_H}
 =\frac12k_BT_H\big|_{R=R_H}.
 \label{rm:eq:horizon-Schwarzschild-temperature}
\end{equation}
En conséquence,
\begin{equation}
 E_H=2T_{\rm Sch}S_H,
 \qquad
 T_H^{\rm dS}=2T_{\rm Sch}
 \label{rm:eq:horizon-Schwarzschild-Smarr}
\end{equation}
pour un même rayon. Dans \(R_H=\ell_P\),
\begin{equation}
 E_H=\frac{E_P}{2},
 \qquad
 T_{\rm Sch}S_H=\frac{E_P}{4}.
 \label{rm:eq:horizon-Schwarzschild-Planck}
\end{equation}
La comparaison sépare les deux géométries : elles partagent l’énergie quasi locale, l’aire et l’entropie, et leurs gravités de surface produisent des températures dans un rapport de deux.

Avec \(\beta=(k_BT_{\rm Sch})^{-1}\), l’action euclidienne satisfait
\begin{equation}
 \frac{I_E}{\hbar}
 =\frac{E_H}{k_BT_{\rm Sch}}-\frac{S_H}{k_B}
 =2\pi-\pi=\pi
 \label{rm:eq:horizon-Schwarzschild-action}
\end{equation}
au rayon de Planck. Le même \(\pi\) apparaît comme entropie réduite, action euclidienne réduite, temps du flot hyperbolique et logarithme du poids effectif.

\subsection{Entrelaceur horizon–Perron : compatibilité spectrale entre énergie de frontière et mémoire de retour}

La coïncidence \(\Omega_{\rm eff}=\rho(\mathcal E_\pi)\) fixe une donnée spectrale et laisse ouverte la relation entre les objets complets. Notons \(\mathcal A_H^{(m)}\) l’algèbre des observables d’horizon au niveau nonadique \(m\), \(\omega_H^{(m)}\) son état, et \(\mathcal P_\varphi^{(m)}\) le système de transfert de Perron.

\begin{definition}[Entrelaceur thermodynamique--dynamique]
\label{rm:def:horizon-Perron-intertwiner}
Un entrelaceur admissible est une famille d’applications linéaires unitales complètement positives (UCP)
\begin{equation}
 \mathfrak I_m:\mathcal A_H^{(m)}\longrightarrow
 \mathcal A_\varphi^{(m)}
 \label{rm:eq:horizon-intertwiner}
\end{equation}
qui satisfait :
\begin{enumerate}[label=\roman*)]
 \item préservation de l’état, en plus de l’unité et de la positivité complète déjà intégrées à la condition UCP ;
 \item compatibilité avec le raffinement nonadique ;
 \item transport de l’involution intérieur--extérieur vers \(J_\varphi\mapsto-J_\varphi\) ;
 \item commutation entre le semi-groupe modulaire d’horizon et le semi-groupe de transfert sur le sous-espace déclaré ;
 \item égalité du rayon spectral au temps \(\pi\), donnée par \cref{rm:eq:horizon-Perron-log}.
\end{enumerate}
\end{definition}

La cinquième condition est nécessaire et scalaire ; les quatre premières fournissent le contenu algébrique qu’une coïncidence numérique ne contient pas. Un entrelaceur construit de cette manière identifierait la multiplicité thermodynamique effective à une dynamique d’expansion et de contraction tout en conservant l’orientation, l’état et le raffinement.

\begin{proofstatusbox}
Les identités \(E_\Lambda=E_{\rm MS}=T_HS_H=c^4R/(2G)\), la définition \(\mathcal Q_H=\mathscr F_P\,R\), la spécialisation CT108, les deux cellules de demi-action et l’égalité \(S_H/k_B=\log\rho(\ee^{\pi J_\varphi})=\pi\) sont des résultats exacts sous les conventions et la réalisation de Sitter déclarées. L’identification de \(\Omega_{\rm eff}\) à un cardinal microscopique et l’entrelaceur \(\mathfrak I_m\) relèvent d’une interface physique délimitée.
\end{proofstatusbox}


\begin{falsifierbox}
La fermeture géométrique échoue si \(\Lambda_R\ne3/R^2\), si la sphère de rayon \(R\) ne satisfait pas la condition marginale ou si l’énergie intégrée diffère de \(E_{\rm MS}(R)\). La fermeture thermodynamique échoue si la température ou la loi d’aire produisent \(T_HS_H\ne c^4R/(2G)\). La spécialisation CT108 échoue si \(R_{108}\ne\ell_P\), \(T_{108}\ne2\pi t_P\) ou \(E_P\ne k_B\Theta_{\rm clk}\log3\). L’interface de Perron est réfutée par une obstruction de positivité, d’état, d’orientation ou de raffinement, même si la valeur propre \(\ee^\pi\) coïncide scalairement.
\end{falsifierbox}
